\chapter{Cross-Sectional Deep Models}\label{ch:ml:cross-sectional-deep-models}

A statistical factor model extracts three factors from a panel of stock returns and gives each stock a fixed loading on
each. A stock listed yesterday has no loading; a stock whose characteristics have changed keeps the loading of the stock
it was. On a simulated panel whose true betas are functions of ten characteristics, half of their variation nonlinear,
the loadings of principal-component analysis explain 22.5\% of the test months' returns, and their expected returns
correlate at 0.03 with the true ones. Betas written as linear functions of the characteristics (instrumented PCA)
explain 27.6\% and correlate at 0.71; betas written as a small \omterm{def:ml:neural-networks-for-noisy-tabular-data:mlp}{neural network} of them (a \omterm{def:ml:cross-sectional-deep-models:ae}{conditional autoencoder})
explain 30.4\% and correlate at 0.95, against a ceiling of 31.3\%. This chapter builds models that learn from a panel as a
panel: pooled across assets, with inputs that describe the asset rather than name it.

\section{Panels, pooling and permutation invariance}

\begin{definition}[Pooled panel model, permutation-invariant network]\label{def:ml:cross-sectional-deep-models:pooled}
A \emph{pooled panel model}\index{pooled panel model} fits one function $f_\theta(x_{it})$ to every asset $i$ and date $t$
of a panel, so that assets share parameters through their characteristics rather than each having its own. A
\emph{permutation-invariant network}\index{permutation-invariant network} maps a set of vectors to an output that does
not depend on their order, typically $\rho(\sum_i\phi(x_i))$ for networks $\phi$ and $\rho$ (Zaheer et al., 2017).
\end{definition}

Every model of chapters~4 to~7 was pooled: one ridge regression or one network for all 500 stocks. Pooling is what makes a
panel of 500 stocks and 240 months a large dataset; its price is the assumption that stocks with the same
characteristics behave alike, which is exactly what a factor model with characteristic-driven betas says. Permutation
invariance appears as soon as a model needs the whole cross-section as an input: the average of any characteristic over
the stocks, or a characteristic-weighted portfolio's return, does not depend on how the stocks are ordered, and it is
how the models below turn a cross-section of any size into a fixed number of factors.

\section{Embeddings for identifiers and categories}

\begin{definition}[Entity embedding]\label{def:ml:cross-sectional-deep-models:embedding}
An \emph{entity embedding}\index{entity embedding} represents each value $j$ of a categorical variable (an industry, a
venue, a stock identifier) by a vector $e_j\in\mathbb R^m$ learned with the rest of the network, in place of a one-hot
column per category (Guo and Berkhahn, 2016).
\end{definition}

Thirty industries are planted on top of the chapter's panel, each with a monthly premium drawn with a standard deviation
of 0.4\%. A pooled network on the ten characteristics alone explains 0.004\% of the test months' demeaned returns; with a
one-dimensional embedding of the industry it explains 0.15\%, its forecasts'
correlation with the true expected returns rises from 0.52 to 0.75, and the learned embedding correlates at 0.86 with the
planted premia (the industries' raw average returns over the training months correlate at 0.88). One practical detail
decided the result: embeddings initialised with PyTorch's default standard deviation of one never moved far enough from
their random start for the premia to show; initialised near zero (standard deviation 0.01), they learned them.

\section{Learned factor models: instrumented PCA and conditional autoencoders}

The chapter's panel is \texttt{firm.mlsynth.factor\_panel}: 300 stocks, 240 months, ten ranked characteristics, three
factors with monthly premia of 0.6\%, 0.4\% and 0.3\%, and betas that are functions of the characteristics, half of their
cross-sectional variance linear and half nonlinear (squares, products, absolute values). The truth's expected return is
$\mu_{it} = \beta(z_{it})^\top\lambda$.

\begin{definition}[Instrumented principal component analysis]\label{def:ml:cross-sectional-deep-models:ipca}
\emph{Instrumented principal component analysis}\index{instrumented principal component analysis} (IPCA) models
$r_{i,t+1} = z_{it}^\top\Gamma f_{t+1} + e_{i,t+1}$ with $K$ latent factors $f$ and betas $z^\top\Gamma$ linear in the
characteristics (a constant included); $\Gamma$ ($L\times K$) and the factors are estimated by alternating least squares,
the factors month by month as cross-sectional regressions of returns on $Z_t\Gamma$ (Kelly, Pruitt and Su, 2019).
\end{definition}

\begin{definition}[Autoencoder, conditional autoencoder]\label{def:ml:cross-sectional-deep-models:ae}
An \emph{autoencoder}\index{autoencoder} is a network trained to reproduce its input through a lower-dimensional
bottleneck: an encoder maps $x$ to a code $h$, a decoder maps $h$ back to $\hat x$, and the loss is the reconstruction error.
A \emph{conditional autoencoder}\index{conditional autoencoder} reconstructs the cross-section of returns $r_{t+1}$ as
$\beta(z_{it})^\top f_{t+1}$, where the code $f_{t+1}$ ($K$ factors) is a linear map of characteristic-managed
portfolios $Z_t^\top r_{t+1}/n$ and the betas are a network of the characteristics (Gu, Kelly and Xiu, 2021).
\end{definition}

The managed portfolios are the permutation-invariant summary of the cross-section: their number is the number of
characteristics, whatever the number of stocks. The chapter's \omterm{def:ml:cross-sectional-deep-models:ae}{autoencoder} has a beta network with two hidden layers (32
and 16 units) and a linear factor map; it is trained on months 0 to 159 with \omterm{def:ml:trees-and-boosting:early}{early stopping} on months 160 to 179; IPCA
and PCA are fitted on months 0 to 179; all three are scored on months 180 to 239.

\begin{proposition}[Why IPCA is a regression]\label{prop:ml:xsnet:ipca}
For fixed $\Gamma$, the least-squares factors of month $t$ are $\hat f_{t+1} = (\Gamma^\top Z_t^\top Z_t\Gamma)^{-1}\Gamma^\top
Z_t^\top r_{t+1}$; for fixed factors, $\operatorname{vec}(\Gamma^\top)$ solves the normal equations
$\sum_t(Z_t^\top Z_t\otimes f_{t+1}f_{t+1}^\top)\operatorname{vec}(\Gamma^\top) = \sum_t(Z_t^\top r_{t+1})\otimes f_{t+1}$.
\end{proposition}

\begin{proof}
The first is ordinary least squares of $r_{t+1}$ on the columns of $Z_t\Gamma$. For the second, $z_{it}^\top\Gamma f =
(z_{it}\otimes f)^\top\operatorname{vec}(\Gamma^\top)$, a linear model in $\operatorname{vec}(\Gamma^\top)$ whose normal
equations sum $(z\otimes f)(z\otimes f)^\top = zz^\top\otimes ff^\top$ over assets and months.
\end{proof}

\section{Evaluating a learned factor model}

Two R-squared measure a factor model (Kelly, Pruitt and Su; Gu, Kelly and Xiu). The \emph{total} R-squared uses the
month's realised factors: how much of the cross-section's variation the betas can explain once the factors are known.
The \emph{predictive} R-squared uses the average factor of the training months, $\beta^\top\hat\lambda$: an expected-return
forecast. Both are measured against zero.

\begin{table}[htbp]
\centering\small
\begin{tabular}{@{}lccccc@{}}
\toprule
 & \multicolumn{3}{c}{total $R^2$} & \multicolumn{2}{c}{3 factors: corr.\ of $\hat\mu$ with}\\
\cmidrule(lr){2-4}\cmidrule(l){5-6}
model & 1 factor & 3 factors & 5 factors & the truth $\mu$ & $\mu$'s nonlinear part\\
\midrule
PCA (static betas) & 22.0\% & 22.5\% & 23.0\% & 0.03 & 0.00\\
IPCA (linear betas) & 25.1\% & 27.6\% & 28.1\% & 0.71 & $-0.03$\\
\omterm{def:ml:cross-sectional-deep-models:ae}{conditional autoencoder} & 26.9\% & 30.4\% & 30.8\% & 0.95 & 0.59\\
\midrule
the truth & \multicolumn{3}{c}{31.3\%} & 1 & \\
\bottomrule
\end{tabular}
\caption{Three factor models on the sixty test months of the conditional factor panel. Total R-squared uses realised
factors; the last two columns correlate each model's expected returns $\beta^\top\hat\lambda$ with the true $\mu$ and with
the part of $\mu$ that comes from the betas' nonlinear terms. \omterm{def:ml:cross-sectional-deep-models:ae}{Autoencoder}: seed~1; seeds 2 and 3 give 30.5\% and 30.6\%
total and correlations of 0.95 and 0.96. Data: \texttt{ml\_xsnet}.}\label{tab:ml:xsnet:table}
\end{table}

\Cref{tab:ml:xsnet:table} has three findings. Static betas are the wrong model for a panel whose characteristics move:
PCA's expected returns carry no cross-sectional information (correlation 0.03), although its total R-squared, driven by
the market factor, is respectable. Linear betas recover the linear part of the truth and none of the nonlinear part
($-0.03$). The \omterm{def:ml:cross-sectional-deep-models:ae}{autoencoder} recovers both, most of the way (0.95 overall, 0.59 of the nonlinear part), and its total
R-squared is within a point of the truth's with three factors. The predictive R-squared, 1.01\% for PCA, 1.11\% for IPCA and
1.23\% for the \omterm{def:ml:cross-sectional-deep-models:ae}{autoencoder} against 1.09\% for the truth itself, ranks the models the same way but should not be read
alone: over sixty months it is dominated by the level of the premia and by how the test window's factor returns happened
to average, which is why the truth can be beaten on it.

\begin{omfigure}
\begin{tikzpicture}
\begin{axis}[width=10.5cm,height=5.4cm,xlabel={number of factors $K$},ylabel={total $R^2$ out of sample (\%)},
  tick label style={font=\small},label style={font=\small},xmin=0.5,xmax=5.5,ymin=20,ymax=33,xtick={1,3,5},grid=major,
  grid style={gray!25},legend columns=4,legend style={font=\footnotesize,at={(0.5,-0.3)},anchor=north,draw=none,
  fill=none,/tikz/every even column/.append style={column sep=6pt}},table/col sep=comma]
\addplot[gray,very thick,mark=triangle*] table[x=K,y=pca]{figdata/ml/09-cross-sectional-deep-models/total.csv};
\addplot[omThm,very thick,mark=square*] table[x=K,y=ipca]{figdata/ml/09-cross-sectional-deep-models/total.csv};
\addplot[omDef,very thick,mark=*] table[x=K,y=cae]{figdata/ml/09-cross-sectional-deep-models/total.csv};
\addplot[black!70,thick,dashed] table[x=K,y=truth]{figdata/ml/09-cross-sectional-deep-models/total.csv};
\legend{PCA,IPCA,conditional autoencoder,truth (3 factors)}
\end{axis}
\end{tikzpicture}

\omcaption{Total R-squared of the test months against the number of factors. The truth has three factors; the models with
more have more room to fit, not more truth to find. Data: \texttt{ml\_xsnet.pca}, \texttt{ipca}, \texttt{cae}.}\label{fig:ml:xsnet:total}
\end{omfigure}

\begin{method}[Fitting a learned factor model]\label{met:ml:xsnet:method}
\begin{enumerate}
\item Rank the characteristics by date and add a constant; keep the return of month $t+1$ with the characteristics of
month $t$.
\item Fit IPCA first: it is fast, has no seeds, and tells how much a linear-beta model explains.
\item Fit the \omterm{def:ml:cross-sectional-deep-models:ae}{conditional autoencoder} with \omterm{def:ml:trees-and-boosting:early}{early stopping} on held-out months and several seeds; compare with IPCA on
total R-squared and on the correlation of expected returns with a benchmark forecast.
\item Choose the number of factors on held-out months; report both R-squared and never the predictive one alone.
\end{enumerate}
\end{method}

\section{Tutorial: factors that learn}\label{tut:ml:cross-sectional-deep-models:tut}

\begin{tutorial}
\textbf{Goal.} Fit PCA, IPCA and a \omterm{def:ml:cross-sectional-deep-models:ae}{conditional autoencoder} with one, three and five factors on a panel with
characteristic-driven betas; measure how much of the planted nonlinearity each recovers; add industries through an
embedding. \textbf{End state:} \cref{tab:ml:xsnet:table}, \cref{fig:ml:xsnet:total}.
\begin{tutsteps}
\item \textbf{IPCA} by alternating least squares (\cref{prop:ml:xsnet:ipca}).
\omcode{code/firm/xsnet/firm_xsnet.py}{66}{89}{Instrumented PCA: factors by cross-sectional regression, loadings by the Kronecker normal equations.}\label{lst:ml:xsnet:ipca}
\item \textbf{The \omterm{def:ml:cross-sectional-deep-models:ae}{conditional autoencoder}}: betas from characteristics, factors from managed portfolios.
\omcode{code/firm/xsnet/firm_xsnet.py}{100}{111}{The conditional autoencoder's two networks.}\label{lst:ml:xsnet:ca}
\item \textbf{Run} \texttt{ml\_xsnet.pca(K)}, \texttt{ipca(K)}, \texttt{cae(K, seed)}, \texttt{embedding()} and
\texttt{fig\_xsnet.py}.
\end{tutsteps}
\textbf{What to change next.} Set the planted nonlinear share to zero and check that IPCA catches up with the
\omterm{def:ml:cross-sectional-deep-models:ae}{autoencoder}; give the \omterm{def:ml:cross-sectional-deep-models:ae}{autoencoder}'s factor map the managed portfolios of squared characteristics too.
\end{tutorial}

\section{Build: cross-sectional networks and factor models}\label{bld:ml:cross-sectional-deep-models:xsnet}

\begin{build}
\bfield{Purpose} The firm's learned factor models, comparable on one panel and one pair of R-squared.
\bfield{Interface} \texttt{managed(Z, r)}, \texttt{with\_constant(Z)}, \texttt{PCAModel(K)}, \texttt{IPCA(K, iters)}
with \texttt{fit}, \texttt{factors}, \texttt{betas}, \texttt{score}; \texttt{CondAE(L, K, hidden, seed)} with the same
methods; \texttt{scores(pred\_total, pred\_mu, r)}; \texttt{EmbedMLP}, \texttt{fit\_embed}.
\bfield{Rules} Characteristics of month $t$ with returns of month $t+1$; managed portfolios computed with the same month's
characteristics and returns; one CPU thread, \omterm{def:ml:trees-and-boosting:early}{early stopping} on held-out months.
\bfield{Acceptance tests} \texttt{code/firm/xsnet/tests/}: IPCA recovers a planted linear-beta model's betas (up to rotation)
and its alternating steps never increase the loss; PCA recovers static loadings; the \omterm{def:ml:cross-sectional-deep-models:ae}{autoencoder}'s managed-portfolio
input is invariant to reordering the stocks; an embedding recovers planted category effects.
\bfield{Stretch} IPCA with an intercept (alpha) test; the \omterm{def:ml:cross-sectional-deep-models:ae}{autoencoder} with nonlinear factor maps; the no-arbitrage
criterion of Chen, Pelger and Zhu.
\end{build}

\begin{omsources}
\item B.~Kelly, S.~Pruitt and Y.~Su, ``Characteristics are covariances: a unified model of risk and return'', \emph{Journal
of Financial Economics} 134(3), 2019.
\item S.~Gu, B.~Kelly and D.~Xiu, ``Autoencoder asset pricing models'', \emph{Journal of Econometrics} 222(1), 2021.
\item L.~Chen, M.~Pelger and J.~Zhu, ``Deep learning in asset pricing'', \emph{Management Science} 70(2), 2024.
\item C.~Guo and F.~Berkhahn, ``Entity embeddings of categorical variables'', arXiv, 2016.
\item M.~Zaheer et al., ``Deep sets'', \emph{NeurIPS}, 2017.
\end{omsources}

\section{Exercises}

\begin{exercise}[$\star$]\label{exo:ml:cross-sectional-deep-models:1}
With ten characteristics plus a constant and three factors, how many parameters has IPCA's $\Gamma$? How many has a static
PCA model of 300 stocks with three factors?
\end{exercise}

\begin{exercise}[$\star$]\label{exo:ml:cross-sectional-deep-models:2}
Three stocks have characteristic values 0.5, $-0.5$, 1.0 and returns 2\%, $-1\%$, 3\% in a month. What is the
characteristic's managed-portfolio return, and does it depend on the order of the stocks?
\end{exercise}

\begin{exercise}[$\star$]\label{exo:ml:cross-sectional-deep-models:3}
Thirty industries, a one-hot encoding into a hidden layer of 32 units against a one-dimensional embedding feeding the same
layer: how many parameters does each use for the industry?
\end{exercise}

\begin{exercise}[$\star\star$]\label{exo:ml:cross-sectional-deep-models:4}
Why does PCA's expected return carry no cross-sectional information on this panel, although its total R-squared is 22.5\%?
\end{exercise}

\begin{exercise}[$\star\star$]\label{exo:ml:cross-sectional-deep-models:5}
Why can a model beat the truth on predictive R-squared over sixty months, and what should be reported instead?
\end{exercise}

\begin{exercise}[$\star\star$]\label{exo:ml:cross-sectional-deep-models:6}
\emph{Find the flaw.} ``Our \omterm{def:ml:cross-sectional-deep-models:ae}{conditional autoencoder} with ten factors explains 40\% of returns in sample against IPCA's 30\%,
so it captures more of the cross-section.''
\end{exercise}

\begin{exercise}[$\star\star\star$]\label{exo:ml:cross-sectional-deep-models:7}
\emph{Coding.} Rerun \texttt{ml\_xsnet.embedding} with embeddings initialised at PyTorch's default (standard deviation one).
Report the embedding's correlation with the premia and explain.
\end{exercise}

\begin{exercise}[$\star\star\star$]\label{exo:ml:cross-sectional-deep-models:8}
Show that IPCA with $\Gamma$ equal to the identity and $K = L$ reduces to the regression of each month's returns on its
characteristics (Fama and MacBeth's first step, Book~4, chapter~16).
\end{exercise}

\section{Problem: Factors That Learn}

\begin{problem}[{Weekend problem --- betas that move with characteristics}]\label{pb:ml:cross-sectional-deep-models:1}
The chapter's panel: 300 stocks, 240 months, ten characteristics, three factors, betas half nonlinear.

\medskip\noindent\textbf{Part I --- The truth.}
\begin{enumerate}
\item What total and predictive R-squared does the truth score on the test months?
\item Why is the truth's total R-squared far from one?
\item What are the premia, and what do they do to predictive R-squared?
\item What does ``half nonlinear'' mean for a model with linear betas?
\end{enumerate}

\medskip\noindent\textbf{Part II --- Three models.}
\begin{enumerate}[resume]
\item Give each model's total R-squared with one, three and five factors.
\item How well do each model's expected returns correlate with the truth and its nonlinear part?
\item How much do seeds move the \omterm{def:ml:cross-sectional-deep-models:ae}{autoencoder}?
\item Why does IPCA recover none of the nonlinear part?
\end{enumerate}

\medskip\noindent\textbf{Part III --- Pooling and embeddings.}
\begin{enumerate}[resume]
\item Why is pooling the right assumption for this panel?
\item What does the industry embedding add, and how close does it get to the raw industry averages?
\item What decided whether the embedding learned at all?
\item Where does permutation invariance appear in the \omterm{def:ml:cross-sectional-deep-models:ae}{autoencoder}?
\end{enumerate}

\medskip\noindent\textbf{Part IV --- The verdict.}
\begin{enumerate}[resume]
\item State the \emph{named result}: out-of-sample total and predictive R-squared of PCA, IPCA and the \omterm{def:ml:cross-sectional-deep-models:ae}{autoencoder} at three
factors, and the share of the planted nonlinearity each recovers.
\item Which model would you use to construct a factor-neutral portfolio, and why?
\item How would you choose the number of factors?
\item What changes with real characteristics that are noisy and missing?
\item What would a stock listed yesterday get from each model?
\item What does the \omterm{def:ml:cross-sectional-deep-models:ae}{autoencoder} cost that IPCA does not?
\item How would you test that the \omterm{def:ml:cross-sectional-deep-models:ae}{autoencoder}'s extra fit is not \omterm{def:ml:why-financial-machine-learning-is-different:generalisation}{overfitting}?
\item In one sentence: what does a learned factor model learn?
\end{enumerate}
\end{problem}

\section{Interview questions}

\begin{interviewq}[$\star$ \iqroles{researcher}]\label{iq:ml:cross-sectional-deep-models:1}
What problem does instrumented PCA solve that PCA on returns does not?
\end{interviewq}

\begin{interviewq}[$\star\star$ \iqroles{researcher, mle}]\label{iq:ml:cross-sectional-deep-models:2}
Explain the \omterm{def:ml:cross-sectional-deep-models:ae}{conditional autoencoder}'s two networks and what goes into each.
\end{interviewq}

\begin{interviewq}[$\star\star$ \iqroles{mle}]\label{iq:ml:cross-sectional-deep-models:3}
When would you use an \omterm{def:ml:cross-sectional-deep-models:embedding}{entity embedding} for a stock identifier, and what can go wrong?
\end{interviewq}

\begin{interviewq}[$\star\star$ \iqroles{researcher}]\label{iq:ml:cross-sectional-deep-models:4}
Total R-squared or predictive R-squared: which do you trust to compare factor models, and why?
\end{interviewq}

\begin{interviewq}[$\star\star$ \iqroles{researcher, mle}]\label{iq:ml:cross-sectional-deep-models:5}
How would you feed a model information about the whole cross-section at a date, for any number of stocks?
\end{interviewq}

\begin{interviewq}[$\star\star\star$ \iqroles{researcher}]\label{iq:ml:cross-sectional-deep-models:6}
Derive IPCA's alternating least-squares steps.
\end{interviewq}
